% ============================================================================
% 2.1 Contexto y motivación
% ============================================================================
\section{Contexto y motivación}

La gesti\'on de inversiones requiere asignar capital entre m\'ultiples activos considerando riesgo, horizonte y condiciones de mercado. FinPUC es una plataforma ficticia de recomendaci\'on de inversiones para clientes del mercado accionario estadounidense clasificados por perfil de riesgo. El sistema genera peri\'odicamente portafolios optimizados buscando maximizar el retorno esperado y las utilidades por comisiones, sujeto a restricciones de riesgo y tiempo computacional.

% ============================================================================
% 2.2 Ciclo operacional
% ============================================================================
\section{Ciclo operacional}

El sistema opera en dos escalas temporales. A nivel de optimizaci\'on, los
pesos del portafolio se recalculan cada semestre (126 d\'ias h\'abiles bursátil) mediante
rebalanceo: se reestiman los par\'ametros $\mu$ y $\Sigma$ con la ventana
hist\'orica acumulada hasta ese punto y se resuelve nuevamente el problema de
optimizaci\'on.

A nivel de simulaci\'on del cliente (Monte Carlo P4), el sistema eval\'ua
semanalmente durante 260 semanas (5 a\~nos) el comportamiento del
inversionista. Cada semana, el cliente puede aceptar la recomendaci\'on vigente
con probabilidad $P_2(x)$ o rechazarla. Si acepta, se ejecuta la
transacci\'on y la empresa cobra una comisi\'on proporcional $k\% = 1\%$ sobre
el monto transado. Si la p\'erdida acumulada del cliente supera la tolerancia
de su perfil, existe una probabilidad $P_1(x)$ de que retire sus fondos de la
plataforma.

% ============================================================================
% 2.3 Perfiles de riesgo
% ============================================================================
\section{Perfiles de riesgo}

Se consideran cinco perfiles de riesgo definidos por una volatilidad anual máxima de 8\%, 12\%, 18\%, 28\% y 40\%, correspondientes a perfiles muy conservador, conservador, neutro, arriesgado y muy arriesgado, respectivamente.

% ============================================================================
% 2.4 Formulación del problema
% ============================================================================
\section{Formulaci\'on del problema}

\subsection{Decisi\'on principal}

Para cada perfil de riesgo y cada punto de rebalanceo, el sistema debe
determinar el vector de pesos $w \in \mathbb{R}^n$ que asigna una fracci\'on
del capital a cada una de las $n = 601$ acciones del universo.

\subsection{Objetivos}

El problema busca maximizar simultáneamente el retorno esperado del cliente ($\mu^\top w$) y las utilidades por comisiones generadas por las transacciones aceptadas.
    
\subsection{Restricciones}
\begin{itemize}
    \item \textbf{No venta en corto:} $w_i \geq 0$ para todo $i$.
    \item\textbf{Presupuesto completo:} $\sum_{i=1}^{n} w_i = 1$.
    \item\textbf{Concentraci\'on m\'axima:} $w_i \leq w_{\max} = 0.20$ para todo $i$.
    \item\textbf{Riesgo por perfil:} $w^\top \Sigma\, w \leq \sigma_{\max}^2$, donde la cota depende del perfil de riesgo.
    \item\textbf{Tiempo de c\'omputo:} el sistema debe generar la recomendaci\'on en menos de 10 segundos.
\end{itemize}


% ============================================================================
% 2.5 Datos disponibles
% ============================================================================
\section{Datos disponibles}

\subsection{Fuente y estructura}

Los datos corresponden a registros hist\'oricos del mercado accionario
estadounidense. El dataset original comprende 1\,406 archivos CSV, uno por ticker, con precios diarios que cubren aproximadamente desde el a\~no 2000
hasta marzo de 2026. Cada archivo incluye fecha, precio de apertura y cierre,
m\'aximo y m\'inimo del d\'ia, volumen de transacciones, dividendos y factores
de \textit{split}. Adicionalmente, se dispone de un archivo de metadatos con
informaci\'on fundamental de cada empresa: sector, industria, tipo de
instrumento y capitalizaci\'on de mercado.

El retorno diario total de cada acci\'on se calcula incorporando dividendos:
\begin{equation}
r_{i,t}=\frac{P_{i,t}-P_{i,t-1}+D_{i,t}}{P_{i,t-1}},
\end{equation}
donde $P_{i,t}$ es el precio de cierre del activo $i$ en el d\'ia $t$ y
$D_{i,t}$ el dividendo pagado en esa fecha.

\subsection{Proceso de filtrado}

A partir del universo crudo de 1\,406 tickers se aplic\'o un proceso de
filtrado secuencial. La Tabla~\ref{tab:filtrado_sin} detalla la cascada del
universo operacional comparable usado en los experimentos, que queda en
601 activos. En paralelo, el Anexo N°2 documenta el universo F5
descriptivo completo de 636 acciones. Esta diferencia no corresponde a dos
universos finales distintos por escenario, sino a dos niveles de an\'alisis:
636 acciones para caracterizaci\'on descriptiva y 601 acciones para
comparaci\'on experimental con y sin pandemia bajo el mismo periodo futuro.

\begin{table*}[h]
\centering
\caption{Cascada de filtrado aplicada al universo sin pandemia.}
\label{tab:filtrado_sin}
\small
\begin{tabular}{@{}p{0.07\textwidth}p{0.67\textwidth}rr@{}}
\toprule
Filtro & Criterio operacional & Tickers & Retenci\'on F0 \\
\midrule
F0 & Archivos CSV encontrados & 1.406 & 100,0\% \\
F1 & Precio $\geq$ USD 5; historia $\geq$ 5 a\~nos; cap. $\geq$ USD 300M o sin dato & 827 & 58,8\% \\
F2 & Excluye \emph{Shell Companies}, sector desconocido y no \texttt{EQUITY} & 824 & 58,6\% \\
F3 & Capitalizaci\'on burs\'atil $\geq$ USD 2B & 659 & 46,9\% \\
F4 & Historia fuera de pandemia $\geq$ 10 a\~nos, aproximada por 2.500 filas & 606 & 43,1\% \\
F5 & Volatilidad anualizada fuera de pandemia entre 5\% y 100\% & 601 & 42,7\% \\
\bottomrule
\end{tabular}
\end{table*}

\subsection{Justificaci\'on del filtrado}

El proceso de filtrado busca construir un universo de activos estable e invertible para la optimizaci\'on de portafolios. F1 elimina acciones con precios muy bajos, historia insuficiente o capitalizaci\'on reducida, disminuyendo estimaciones poco robustas. F2 restringe el an\'alisis a instrumentos \texttt{EQUITY} y excluye \textit{Shell Companies} y sectores desconocidos para mantener comparabilidad entre activos.

Luego, F3 exige una capitalizaci\'on m\'inima de USD 2B para reducir problemas de liquidez y volatilidad extrema. F4 requiere al menos diez a\~nos de historia fuera del periodo pandemia para estimar par\'ametros con mayor estabilidad. Finalmente, F5 elimina activos con volatilidad anualizada excesivamente alta o baja, asociados potencialmente a errores de datos o comportamientos incompatibles con un portafolio accionario tradicional.


\subsection{Dise\~no experimental: escenarios con y sin pandemia}

La evaluaci\'on del sistema se realiza bajo dos escenarios que comparten el
mismo per\'iodo futuro de evaluaci\'on (2 a\~nos fuera de muestra) pero
difieren en la ventana de entrenamiento. El escenario ``sin pandemia'' excluye
las observaciones de 2020 - 2022, estimando $\mu$ y $\Sigma$ con un mercado
relativamente normal. El escenario ``con pandemia'' las incluye, incorporando
correlaciones de estr\'es y volatilidad extrema. Cualquier diferencia en
resultados entre ambos escenarios se interpreta como efecto de la
informaci\'on usada para calibrar, no como diferencias de mercado futuro.

\subsection{Riesgos y limitaciones de los datos}

\begin{itemize}
    \item Posible sesgo de supervivencia por exigir historia larga y alta capitalización.
    \item Sobrerrepresentación del sector financiero en el universo final.
    \item Tratamiento alternativo del periodo 2020--2022 mediante escenarios con y sin pandemia.
    \item Diferencias entre retornos usados en filtrado y análisis final.
    \item Pérdida de observaciones por uso de \texttt{join="inner"}.
\end{itemize}
